% Honest Durability for Graph ANN — arXiv-style preprint.
% Compiles with a basic TeX install:  pdflatex honest-durability.tex  (x2 for refs)
% Self-contained: standard `article` class, no external .cls or .bib.
% NOTE: citation details (arXiv IDs / venues) are from memory and should be
% verified before submission (web was unavailable while drafting).
\documentclass[11pt]{article}

\usepackage[margin=1in]{geometry}
\usepackage{amsmath}
\usepackage{booktabs}
\usepackage{graphicx}
\usepackage{xcolor}
\usepackage{listings}
\usepackage{microtype}
\usepackage[hidelinks]{hyperref}

\lstset{basicstyle=\ttfamily\small,breaklines=true,columns=fullflexible}
\newcommand{\fullfsync}{\texttt{F\_FULLFSYNC}}
\newcommand{\fsync}{\texttt{fsync}}

\title{\textbf{Honest Durability for Graph-Based\\ Approximate Nearest Neighbor Indexes}}
\author{Habib Rahman\\ \small \texttt{<affiliation / email>}}
\date{\today\\[2pt]\small Preprint --- work in progress}

\begin{document}
\maketitle

\begin{abstract}
Research on approximate nearest neighbor (ANN) search is evaluated almost
entirely on two axes: recall and queries per second. Durability---whether an
acknowledged write survives a crash or power loss---is largely absent from ANN
systems and from every mainstream ANN benchmark. We argue this omission both
hides a real cost and leaves an under-explored design space. Using a from-scratch
C++20 vector database with a write-ahead-logged, segmented HNSW store, we make
four contributions. (1) A \emph{durability-aware measurement protocol} for ANN
indexes: a tmpfs guard, an \fsync{}-floor baseline, and mandatory
\fsync{}-mode labeling, because the default \fsync{} on some platforms does not
flush the drive cache and thus reports optimistic ``durable'' numbers. (2) The
first ANN-isolated \emph{durability tax}: on our testbed, honest per-write
durability (\fullfsync{}) reduces the device's durable-sync ceiling by
$\sim$130$\times$ (41{,}101 to 316 syncs/s) and single-insert throughput from
442 to 110 inserts/s. (3) \emph{Group commit} for a graph-ANN WAL---one \fsync{}
per batch---recovers $4.4\times$ insert throughput under honest durability and
collapses the p99 tail, while remaining crash-consistent (validated by a
fork-plus-kill test). (4) A \emph{recovery} analysis showing snapshot-load
(sealed segments) recovers a graph index $\sim$6$\times$ faster than WAL replay
(35\,ms vs.\ 225\,ms for 1{,}500 vectors), with a locatable crossover. All code,
harnesses, and results are open and reproducible.
\end{abstract}

\section{Introduction}
Graph-based ANN indexes---HNSW~\cite{hnsw} and its descendants---dominate
in-memory vector search, and vector databases built on them now back
retrieval-augmented generation and semantic search at scale. The community
evaluates them on \emph{recall@$k$} and \emph{throughput} (QPS), summarized by the
recall-vs-QPS Pareto frontier popularized by ANN-Benchmarks~\cite{annbench} and
the Big ANN challenges~\cite{bigann}. What these evaluations almost never measure
is \emph{durability}: does a write the system acknowledged survive a crash?

This is not a niche concern. Vector databases are increasingly operational
stores that ingest and update continuously, and an operational store that loses
acknowledged writes is not a database. Yet mainstream ANN libraries scope
durable storage out entirely (Faiss offers manual serialization, no
WAL~\cite{faiss}; hnswlib has soft-delete tombstones and no crash
consistency~\cite{hnsw}), and no widely used ANN benchmark
(ANN-Benchmarks~\cite{annbench}, Big ANN~\cite{bigann},
VectorDBBench~\cite{vectordbbench}) includes a crash-recovery, \fsync{}, or
durability case. The closest systems either integrate ANN into an operational DB
and \emph{inherit} durability from the host engine---delegating rather than
analyzing it~\cite{cosmos}---or, like LSM-VEC~\cite{lsmvec}, build an LSM-tree
graph index yet report no \fsync{} cost, crash consistency, or recovery time.

We take the opposite stance: build durability into a \emph{standalone} graph-ANN
index, measure it honestly, and characterize the cost. Two facts make this
worth a paper. First, durable writes are expensive in a way ANN work has never
quantified, and the expense is easy to hide: on macOS (and any stack where
\fsync{} does not force a drive-cache flush) the ``durable'' number is a page-cache
number. Second, once you pay the cost, standard database techniques---group
commit and snapshot-based recovery---apply, but their effect on an ANN index
(where each insert also mutates a proximity graph) has not been reported.

\paragraph{Contributions.}
\begin{itemize}
\item \textbf{C1 (\S\ref{sec:honest}).} A durability-aware measurement protocol
for ANN: tmpfs detection, an \fsync{}-floor baseline, and \fsync{}-mode labeling,
so ANN write/durability numbers are comparable and honest.
\item \textbf{C2 (\S\ref{sec:tax}).} The first ANN-isolated durability-tax
characterization: throughput, latency percentiles, and build cost vs.\ the
strength of the durability guarantee.
\item \textbf{C3 (\S\ref{sec:groupcommit}).} Group commit for a graph-ANN WAL and
its measured throughput/tail-latency benefit, with a crash-consistency test.
\item \textbf{C4 (\S\ref{sec:recovery}).} A recovery analysis---snapshot-load vs.\
WAL-replay---for a standalone graph ANN, including the crossover between them.
\end{itemize}

\section{Background and System}
\label{sec:background}
Our testbed is a from-scratch C++20 vector database. The default storage engine
is a Qdrant/LSM-style \emph{segmented store}: a single mutable segment backed by
a write-ahead log (WAL), plus $N$ immutable \emph{sealed} segments, each carrying
an HNSW snapshot. Writes append to the WAL; when a size or tombstone-ratio
threshold is crossed the mutable segment is sealed (its HNSW graph and vectors
serialized), and sealed segments are compacted in the background. A cross-segment
map resolves the current location of each key (MVCC by sequence number).

\paragraph{Durability model.} Every WAL record is framed with a header (magic,
version, op, sequence, field lengths, dimensions) and a CRC-32 covering the
header and payload, and---by default---is \fsync{}'d before the call returns.
Snapshot and manifest writes are published atomically (temp file $\to$ \fsync{}
$\to$ rename $\to$ directory \fsync{}). Recovery of a sealed segment loads its
HNSW snapshot directly (memcpy of neighbor lists and precomputed distances) rather
than replaying the WAL through the graph builder; the mutable segment replays its
(small) WAL. WAL replay stops at the first record failing the magic/version/CRC
check, tolerating a torn tail.

\paragraph{Index.} Search uses HNSW with the Malkov--Yashunin diversifying
neighbor-selection heuristic~\cite{hnsw}; distance kernels are SIMD (NEON/AVX+FMA)
with a devirtualized hot path. After fixing an initially na\"ive
neighbor-selection rule, recall@10 on 30k clustered 128-d vectors rises from a
capped $0.36$ to $0.97$ at $\mathrm{ef}=10$ and $1.0$ by $\mathrm{ef}=100$
(Table~\ref{tab:recall}), i.e.\ the index is competitive before we layer
durability on top.

\section{Measuring Durability Honestly}
\label{sec:honest}
Durability numbers are worthless without stating \emph{what} was made durable.
Two traps recur.

\paragraph{The tmpfs trap.} Benchmarks frequently run under \texttt{/tmp}, which
on many Linux distributions is a RAM-backed \texttt{tmpfs} where \fsync{} is a
no-op. A ``durable'' insert benchmark on tmpfs measures the CPU path only. Our
harness \texttt{statfs}-checks the target and \emph{refuses} to report write
numbers on tmpfs.

\paragraph{The \fsync{}-mode trap.} On macOS, \fsync(2) does not flush the drive's
write cache; true power-loss durability requires \texttt{fcntl(F\_FULLFSYNC)},
which is far slower. A number labeled ``durable'' under plain \fsync{} on macOS
overstates durability. We therefore (a) expose an \fullfsync{} toggle and (b)
report an \emph{\fsync{} floor}: a tight \texttt{write()+sync} loop giving the
device's durable-sync ceiling, so index overhead can be separated from the
device's flush cost. Table~\ref{tab:floor} shows the floor differs by
$\sim$130$\times$ between the two modes on our testbed---the single most important
number to disclose.

\begin{table}[t]
\centering
\caption{\fsync{} floor on the testbed (Apple M4 Pro, APFS/NVMe): durable syncs/s
in a tight \texttt{write()+sync} loop. Plain \fsync{} does not flush the drive
cache; \fullfsync{} does.}
\label{tab:floor}
\begin{tabular}{lr}
\toprule
Mode & Durable syncs/s \\
\midrule
Plain \fsync{} (page cache) & 41{,}101 \\
\fullfsync{} (drive-cache flush) & 316 \\
\bottomrule
\end{tabular}
\end{table}

\section{The Durability Tax}
\label{sec:tax}
We define the \emph{durability tax} as the throughput and latency an index gives
up to guarantee that acknowledged writes survive. We measure it by inserting one
record at a time (each \fsync{}'d before returning) into the segmented store and
varying only the \fsync{} mode. Table~\ref{tab:tax} reports the result on
synthetic data ($d{=}64$); the index work per insert (a real HNSW insertion with
$\mathrm{ef}_{\mathrm{construction}}{=}200$) is held constant, so the delta isolates the
durability cost.

\begin{table}[t]
\centering
\caption{Durability tax: single-insert throughput and latency vs.\ \fsync{} mode
(Apple M4 Pro, $d{=}64$, per-write durability). Honest durability cuts throughput
$4\times$ and inflates the p99 tail $3\times$.}
\label{tab:tax}
\begin{tabular}{lrrr}
\toprule
Mode & Inserts/s & p50 (\textmu s) & p99 (\textmu s) \\
\midrule
Plain \fsync{} & 442 & 1{,}753 & 14{,}351 \\
\fullfsync{} & 110 & 7{,}082 & 43{,}045 \\
\bottomrule
\end{tabular}
\end{table}

Two readouts. First, honest durability is not free: throughput falls $4\times$
and the p99 tail triples, because each insert now waits on a real drive flush.
Second---and this is why the \fsync{} floor matters---under plain \fsync{} the
bottleneck is the graph insertion, not the sync, so a system that stops at plain
\fsync{} both understates its true durability cost \emph{and} understates how much
a batching optimization can help. That is the subject of \S\ref{sec:groupcommit}.

\section{Group Commit for a Graph-ANN WAL}
\label{sec:groupcommit}
Group commit is a classic database technique~\cite{aether}: batch several log
records and issue a single \fsync{} for the group, decoupling per-operation commit
latency from flush latency. We add a deferred-sync mode to the mutable segment:
during a batch, WAL appends skip the per-record \fsync{}, and a single
\texttt{commitDeferredSync()} at the end \fsync{}s once for the whole batch. The
public batch-insert path uses it; single inserts keep per-write durability.

\begin{table}[t]
\centering
\caption{Group commit vs.\ per-write \fsync{}, honest durability
(\fullfsync{}, Apple M4 Pro, $d{=}64$). One \fsync{} per batch recovers
$4.4\times$ throughput.}
\label{tab:gc}
\begin{tabular}{lrr}
\toprule
Path & Inserts/s & Speedup \\
\midrule
Per-write \fsync{} (\fullfsync{}) & 110 & 1.0$\times$ \\
Group commit (\fullfsync{}) & 483 & 4.4$\times$ \\
\bottomrule
\end{tabular}
\end{table}

Under honest durability, group commit lifts insert throughput from 110 to 483
inserts/s ($4.4\times$; Table~\ref{tab:gc}) and collapses the p99 tail (43\,ms
per-write $\to$ amortized across the batch), because the batch pays for one flush
rather than $N$.

\paragraph{Crash consistency.} Group commit is only correct if the batch's single
\fsync{} makes the \emph{whole} batch durable before the call returns. We validate
this with a fault-injection test: a child process performs a group-committed batch
of $N$ inserts and then \texttt{\_exit(0)} \emph{without} a clean shutdown
(simulating a crash immediately after commit); the parent reopens the store and
asserts all $N$ rows recover. It passes. The WAL is protected by a CRC-32 over the
header and payload (an earlier version used a payload-only FNV-1a mislabeled as
``crc32''---replaced here), and replay truncates at the first bad record.

\paragraph{Recall staleness (a note toward a stronger guarantee).} Group commit
opens a window in which recently inserted-but-not-yet-\fsync{}'d vectors are
visible to queries but not yet durable. In an ANN setting this window can be
expressed in \emph{recall} units: bounding the number of un-synced records bounds
the recall a query can lose if a crash discards them. Making this a provable
``a crash never costs more than $\varepsilon$ recall@$k$'' guarantee---a
recall-aware commit scheduler---is a promising direction we leave to future work.

\section{Fast, Crash-Consistent Recovery}
\label{sec:recovery}
Recovery time is a first-class operational metric that ANN work has not reported
for a standalone index. Our store has two recovery paths: a \emph{sealed} segment
loads its HNSW snapshot directly (no graph rebuild), while a \emph{mutable}
segment replays its WAL, re-inserting each record through the graph builder.
Table~\ref{tab:recovery} shows the sealed path is $\sim$6$\times$ faster at
$N{=}1{,}500$ and the gap widens with $N$ (WAL replay is $O(N)$ graph insertions;
snapshot load is a linear read plus memcpy).

\begin{table}[t]
\centering
\caption{Cold-open recovery time, sealed (snapshot load) vs.\ mutable (WAL replay),
Apple M4 Pro, $d{=}64$.}
\label{tab:recovery}
\begin{tabular}{lrr}
\toprule
$N$ & Sealed (ms) & WAL replay (ms) \\
\midrule
1{,}500 & 35 & 225 \\
\bottomrule
\end{tabular}
\end{table}

The practical consequence is a \emph{crossover}: an unsealed WAL that grows too
large blows the recovery SLA, which is the operational justification for periodic
sealing. Fully charting recovery-time-vs-$N$ and this crossover---and validating
consistency under true power loss (Linux \texttt{dm-log-writes})---is the main
remaining measurement (\S\ref{sec:threats}).

\section{Index Competitiveness}
\label{sec:index}
So that the durability results are not measured on a broken index, we confirm the
HNSW is competitive. Table~\ref{tab:recall} traces recall@10 vs.\ QPS on 30k
clustered 128-d vectors; recall is measured against exact ground truth from a flat
index using the identical distance kernel (tie-aware). Recall climbs
monotonically with $\mathrm{ef}$ to $1.0$; after hot-path optimization (a
versioned visited set, FMA distance kernels, devirtualized dispatch, and a pooled
graph arena that cut peak index memory from 2\,GB to 38\,MiB) single-thread QPS
reaches $\sim$97k at $\mathrm{ef}{=}10$.

\begin{table}[t]
\centering
\caption{Recall@10 vs.\ single-thread QPS (Apple M4 Pro, $d{=}128$, $n{=}30$k,
$M{=}16$, $\mathrm{ef}_{c}{=}200$), vs.\ exact ground truth.}
\label{tab:recall}
\begin{tabular}{rrr}
\toprule
$\mathrm{ef}$ & recall@10 & QPS \\
\midrule
10  & 0.97 & 96{,}805 \\
48  & 1.00 & 44{,}041 \\
100 & 1.00 & 31{,}511 \\
\bottomrule
\end{tabular}
\end{table}

\section{Related Work}
\label{sec:related}
\textbf{Graph ANN and updates.} HNSW~\cite{hnsw} and DiskANN~\cite{diskann} are the
dominant graph indexes; FreshDiskANN~\cite{freshdiskann} and SPFresh~\cite{spfresh}
add streaming updates (out-of-place merge and in-place rebalancing respectively),
and LSM-VEC~\cite{lsmvec} distributes a proximity graph across LSM levels. None
report an \fsync{}-tax, crash-consistency, or recovery-time analysis; LSM-VEC in
particular is silent on \fsync{}/recovery despite an LSM design.
\textbf{Disk-resident ANN.} Starling~\cite{starling} optimizes on-disk graph
layout \emph{statically}; SPANN~\cite{spann} splits centroids/posting lists across
RAM/SSD. \textbf{Durability inherited, not analyzed.} Azure Cosmos DB embeds
DiskANN and inherits durability/recovery from the host engine~\cite{cosmos}, so it
does not isolate an ANN-specific recovery cost. \textbf{Benchmarks.}
ANN-Benchmarks~\cite{annbench}, Big ANN~\cite{bigann} (including its in-memory
streaming track), and VectorDBBench~\cite{vectordbbench} measure recall/QPS/cost
but have no crash-recovery or \fsync{} case. \textbf{Durability techniques.} Group
commit dates to early transaction systems and was refined in Aether~\cite{aether};
we apply it to a graph-ANN WAL and characterize its ANN-specific effect.

\section{Threats to Validity and Remaining Work}
\label{sec:threats}
Numbers here are preliminary and measured on a single host (Apple M4 Pro,
APFS/NVMe) using synthetic clustered data; every durability number is labeled by
\fsync{} mode. Because macOS \fsync{} does not flush the drive cache, the honest
figures use \fullfsync{}; a Linux host with \texttt{ext4/btrfs} on NVMe (and
\texttt{dm-log-writes} for true power-loss injection) is needed to generalize.
Remaining work before camera-ready: (i) SIFT1M/GIST1M runs with an hnswlib
recall/QPS parity check; (ii) recovery-time-vs-$N$ curves and the sealed/mutable
crossover on real media; (iii) a concurrent (leader/follower) group committer;
(iv) elevating the recall-staleness note into a proven bound.

\section{Conclusion}
Durability is the axis ANN evaluation forgot. We showed it is measurable, that it
carries a large and previously unreported tax on a graph index, that group commit
recovers most of the throughput while staying crash-consistent, and that
snapshot-based recovery restores a graph index far faster than WAL replay. We hope
the honest-measurement protocol (tmpfs guard, \fsync{}-floor, \fsync{}-mode
labeling) becomes a default for durable ANN systems.

\section*{Reproducibility}
All code, benchmarks, and tests are open. Key entry points:
\texttt{make bench-durability [DURABILITY\_ARGS="--full-fsync --dir /nvme"]}
(Tables~\ref{tab:floor}--\ref{tab:gc},~\ref{tab:recovery}),
\texttt{make bench-ann ANN\_ARGS="--data datasets/sift"} (Table~\ref{tab:recall}),
\texttt{make bench-hnswlib} (baseline), and \texttt{make test} (111 tests, clean
under AddressSanitizer, UndefinedBehaviorSanitizer, and ThreadSanitizer).

\begin{thebibliography}{99}
\small
% NOTE: verify arXiv IDs / venues before submission (web unavailable at drafting).
\bibitem{hnsw} Y.~Malkov and D.~Yashunin. Efficient and robust approximate nearest
neighbor search using Hierarchical Navigable Small World graphs. \emph{IEEE TPAMI}, 2020.
\bibitem{diskann} S.~Subramanya et al. DiskANN: Fast accurate billion-point nearest
neighbor search on a single node. \emph{NeurIPS}, 2019.
\bibitem{freshdiskann} A.~Singh et al. FreshDiskANN: A fast and accurate graph-based
ANN index for streaming similarity search. arXiv:2105.09613, 2021.
\bibitem{spfresh} Y.~Xu et al. SPFresh: Incremental in-place update for
billion-scale vector search. \emph{SOSP}, 2023.
\bibitem{starling} M.~Wang et al. Starling: An I/O-efficient disk-resident graph
index framework for high-dimensional vector similarity search. \emph{SIGMOD}, 2024.
\bibitem{spann} Q.~Chen et al. SPANN: Highly-efficient billion-scale approximate
nearest neighbor search. \emph{NeurIPS}, 2021.
\bibitem{lsmvec} LSM-VEC: A disk-based dynamic vector index over LSM-trees.
arXiv:2505.17152, 2025.
\bibitem{cosmos} Cost-effective, low-latency vector search with Azure Cosmos DB.
arXiv:2505.05885, 2025.
\bibitem{faiss} M.~Douze et al. The Faiss library. arXiv:2401.08281, 2024.
\bibitem{annbench} M.~Aum\"uller, E.~Bernhardsson, A.~Faithfull. ANN-Benchmarks: A
benchmarking tool for approximate nearest neighbor algorithms. \emph{Information
Systems}, 2020.
\bibitem{bigann} H.~Simhadri et al. Results of the NeurIPS'21/'23 Big ANN
(Practical Vector Search) challenges.
\bibitem{vectordbbench} Zilliz. VectorDBBench: A benchmark for vector databases.
\bibitem{aether} R.~Johnson et al. Aether: A scalable approach to logging.
\emph{PVLDB}, 2010.
\end{thebibliography}

\end{document}
